\documentclass[11pt,a4paper]{ctexart}

\usepackage[margin=2.5cm]{geometry}
\usepackage{amsmath,amssymb}
\usepackage{booktabs}
\usepackage{array}
\usepackage{enumitem}
\usepackage{hyperref}

\hypersetup{
  colorlinks=true,
  linkcolor=blue,
  urlcolor=blue,
  citecolor=blue,
  pdftitle={代数几何讨论班 Syllabus}
}

% ===== 可修改的基本信息 =====
\newcommand{\seminartitle}{Seminar on Alegbraic Geometry 2026}
\newcommand{\organizer}{章梓杭}
\newcommand{\seminaremail}{\href{mailto:zhangzh12026@shanghaitech.edu.cn}{zhangzh12026@shanghaitech.edu.cn}}
\newcommand{\seminartime}{待定}
\newcommand{\seminarlocation}{待定}
\newcommand{\seminarterm}{2026 -- 2027 第一学期}
\newcommand{\seminartext}{R. Hartshorne, \emph{Algebraic Geometry}, GTM 52, Springer.}

\title{\textbf{\seminartitle}}
\author{组织者：\organizer \\ \texttt{\seminaremail}}
\date{\seminarterm}

\begin{document}
\maketitle

\subsection*{Seminar Description.}
本讨论班以 Hartshorne《Algebraic Geometry》为主线, 
学习代数几何的基本概念. 我们计划在一学期内共读 
Chapter II.1 -- Chapter II.7, 目标掌握概形的基本性质与工具. 

Hartshorne 作为代数几何的 ``圣经'', 又具有晦涩和习题繁杂的特点, 
因此我们鼓励同学们以该书为主线辅以其他材料进行学习, 并欢迎在讨论班上
分享不同的视角与观点. 此外, 由于 Hartshorne 的正文多引用习题, 
并且有很多重要的命题被写成习题, 本讨论班鼓励同学们讲述一些重要的习题, 
欢迎一切关于习题和定理的讨论. 

讨论班采用轮值主讲和课后讨论的形式, 每周一次, 每次时长 1.5 小时. 
鼓励主讲人在讨论班开始前发布整理好的笔记, 供听众更好理解. 

由于组织者们均为代数几何初学者, 而 Hartshorne 是本新手不友好的书, 
因此我们希望讨论班能够通过共读做到新手友好. 欢迎所有同学来听! 
如您愿意加入讨论班担任主讲, 请联系我们! 

\subsection*{Seminar Information.}

\begin{tabular}{@{}p{3cm}p{11cm}@{}}
\toprule
时间 & \seminartime \\
地点 & \seminarlocation \\
教材 & \seminartext \\
范围 & Chapter II.1 -- Chapter II.7 \\
对象 & 本科生/研究生 \\
预备知识 & 范畴论, 抽象代数, 交换代数 \\
\bottomrule
\end{tabular}

\subsection*{Textbooks \& References.}
\begin{itemize}
    \item R. Hartshorne, \href{https://link.springer.com/book/10.1007/978-1-4757-3849-0}{\emph{Algebraic Geometry}}. 
    \item R. Vakil, \href{https://math.stanford.edu/~vakil/216blog/FOAGjul2724public.pdf}{The Rising Sea: Foundations of Algebraic Geometry}.
    \item D. Eisenbud and J. Harris, \href{https://link.springer.com/book/10.1007/b97680}{\emph{The Geometry of Schemes}}. 
    \item The Stacks project authors, \href{https://stacks.math.columbia.edu/}{The Stacks project}.
\end{itemize}
\hspace{3em} and so on. 


\subsection*{Tentative Schedule.}

\begin{table}[h]
\centering
\begin{tabular}{l l l l} % 四列，内容居中对齐
\toprule[1pt] % 第一条横线
  & Topics & Hartshorne & Speaker \\
\midrule % 第二条横线
1 & Sheaves, stalks, local nature of sheaves, sheafification & II.1 & 章梓杭 \\
2 & $ \mathsf{Sh} $, base sheaves and glueing sheaves, locally ringed space  & II.1 & 章梓杭 \\
3 & Affine schemes, morphisms to affine schemes, glueing schemes & II.2 & 章梓杭 \\
4 & $\operatorname{Proj} S$, functor from $ \mathsf{Var}_{k} $ to $ \mathsf{Sch}_{k} $, functor of points & II.2 &  \\
5 & Reduced, integral, noetherian, finite type, immersion & II.3 &  \\
6 & Algebraic scheme, dimension, fibred product & II.3 &  \\
7 & Separated and proper morphism, valuative criterion & II.4 &  \\
8 & Projective morphisms & II.4 &  \\
9 & Sheaves of modules & II.5 &  \\
10 & Quasi-coherent sheaves & II.5 &  \\
11 & Quasi-coherent sheaves on $ \operatorname{Proj}S $ & II.5 &  \\
12 & Weil divisors & II.6 &  \\
13 & Cartier divisors & II.6 & \\
14 & Projective morphisms, ample invertible sheaves & II.7 &  \\
15 & Linear system, relative $ \mathbf{Proj} $  & II.7 &  \\
 
\bottomrule[1pt] % 第三条横线
\end{tabular}
\end{table}

\end{document}